\banner{Hysteresis, films and the connectivity of the pore network}\label{ch:pg-deep}
\section{Hysteresis of a single pore: the ink-bottle effect}\label{sec:pg-ink}
A real pore is not a uniform tube. It consists of wide \kw{bodies} joined by narrow \kw{throats}, and the two parts control the filling and the emptying in different ways. Consider the idealised pore of Fig.~\ref{fig:pg-ink}, a spherical body of radius $r_b$ connected to the next pore by a cylindrical throat of radius $r_t<r_b$.

\subhd{Drying.} The body is full and the suction rises. The air must enter through the throat, so the meniscus in the throat must be curved enough to balance the suction. From the Young--Laplace law (Equation~\ref{eq:pg-cap}) the body empties at the suction
\begin{equation}\label{eq:pg-hd}\boxed{h_d=\frac{2\gamma\cos\beta}{\rho_w\,g\,r_t}}\end{equation}
At this suction the throat lets the air in and the whole body empties at once, although a body of radius $r_b$ alone would have emptied at a much lower suction. This sudden emptying is a \kw{Haines jump}.

\subhd{Wetting.} The body is empty and the suction falls. The water must fill the body, and the meniscus of largest radius that it must pass is in the body itself; the body therefore fills only when the suction has fallen to
\begin{equation}\label{eq:pg-hw}\boxed{h_w=\frac{2\gamma\cos\beta}{\rho_w\,g\,r_b}<h_d}\end{equation}
Between $h_w$ and $h_d$ the pore can be either full or empty, according to its history. A soil is a population of such pores; the drying (main drainage) curve of the soil is governed by the throats and lies above the wetting (main imbibition) curve, which is governed by the bodies. A reversal of the direction at an intermediate suction follows a \kw{scanning curve} between the two main curves. The ratio $h_d/h_w=r_b/r_t$ is the \kw{aspect ratio} of the pore; for sands it is about 2 and for structured loams it can exceed 5.

\par\noindent\textbf{Notes:} 1. The contact angle of an advancing meniscus is larger than that of a receding one (\kw{contact-angle hysteresis}), which widens the loop further. 2. Air trapped in bodies during wetting is the reason why a soil that is wetted from below rarely reaches the full porosity: the \kw{satiated} water content is typically 0.85 to 0.95 of $\theta_s$. 3. Daisy and the layers of SWAFL use one (drying) retention curve; the error made by ignoring hysteresis is largest in the wet range after rain on a dry soil.

\begin{kfig}{The ink-bottle pore. A body of radius $r_b$ is joined through a throat of radius $r_t$. It empties when the meniscus in the throat can hold the suction $h_d=2\gamma/(\rho_w g r_t)$, and fills again only when the suction has fallen to $h_w=2\gamma/(\rho_w g r_b)$.\label{fig:pg-ink}}
\draw[soil] (-4.2,-2.2) rectangle (4.2,2.2);
\fill[white] (0,0) circle (1.6);\fill[white] (-4.2,-0.35) rectangle (-1.45,0.35);\fill[white] (1.45,-0.35) rectangle (4.2,0.35);
\draw[ln] (0,0) circle (1.6);
\fill[wtr] (0,0) circle (1.55);\fill[khcyan!25] (-4.2,-0.33) rectangle (-1.5,0.33);
\draw[crv] (-1.5,0.33) to[bend left=35] (-1.5,-0.33);
\draw[vec] (0,0)--(1.1,1.1);\node[lab,fill=white,inner sep=1pt] at (0.75,0.35) {$r_b$};
\draw[<->,ln] (2.8,-0.35)--(2.8,0.35);\node[lab,right] at (2.85,0) {$2r_t$};
\node[lab] at (-3.1,0.75) {throat};\node[lab] at (0,-1.95) {body};
\node[lab,align=center] at (0,2.75) {drying: empties at $h_d\propto 1/r_t$ \quad wetting: fills at $h_w\propto 1/r_b$};
\end{kfig}

\exbox{Haines jump in an ink-bottle pore}{A pore of a sandy loam has a body of radius $25\ \mu$m and throats of radius $5\ \mu$m. Take $\gamma=0.0728$ N\,m$^{-1}$, $\beta=0$, $\rho_w=1000$ kg\,m$^{-3}$ and $g=9.81$ m\,s$^{-2}$. Find the suctions at which the pore empties and fills, and the width of the hysteresis loop.}\sol
\givenline{$r_b=25\ \mu\text{m}=25\times10^{-6}$ m ; $r_t=5\ \mu\text{m}=5\times10^{-6}$ m ; $\gamma=0.0728$ N\,m$^{-1}$ ; $\beta=0$ ; $\rho_w=1000$ kg\,m$^{-3}$ ; $g=9.81$ m\,s$^{-2}$}
\sh{Suction at which the pore empties}
\wkt\ during drying the air must enter through the throat, so the suction at which the pore empties is given by Equation~\ref{eq:pg-hd}, with $\cos\beta=1$:
\[h_d=\frac{2\gamma\cos\beta}{\rho_w\,g\,r_t}\ti{i}\]
Substituting the given values in equation \rfi{i},
\[h_d=\frac{2\times0.0728\times1}{1000\times9.81\times5\times10^{-6}}=\frac{0.1456}{0.04905}=2.968\ \text{m}\]
\[\therefore\quad h_d=297\ \text{cm}\Ans\]
\sh{Suction at which the pore fills}
During wetting the meniscus of largest radius lies in the body. Therefore from Equation~\ref{eq:pg-hw},
\[h_w=\frac{2\gamma\cos\beta}{\rho_w\,g\,r_b}=\frac{0.1456}{1000\times9.81\times25\times10^{-6}}=\frac{0.1456}{0.24525}=0.594\ \text{m}=59\ \text{cm}\Ans\]
\sh{Width of the hysteresis loop}
Dividing equation \rfi{i} by the corresponding expression for $h_w$,
\[\frac{h_d}{h_w}=\frac{r_b}{r_t}=\frac{25}{5}=5\Ans\]
Thus between 59 cm and 297 cm the pore may be either full or empty, according to its history. A tensiometer reading of 150 cm in this soil therefore corresponds to two different water contents, the wetter one lying on the drying branch.

\section{Water films and corner flow in the dry range}\label{sec:pg-film}
When the suction is so high that even the smallest throats have emptied, water remains in two forms: in the \kw{corners} of angular pores, held by menisci of small radius, and as \kw{adsorbed films} on the walls, held by the attraction of the solid (van der Waals forces). The film is thin, but it is continuous, and it carries the slow flow that feeds evaporation and the uptake of roots in a dry soil.

\subhd{Thickness of a film.} For a flat wall the energy of the film per unit volume is balanced by the suction. With the Hamaker constant $A_H$ of the solid--water--air system (about $6\times10^{-20}$ J for silicates) the equilibrium thickness is
\begin{equation}\label{eq:pg-film}\boxed{\delta=\Bigl(\frac{A_H}{6\pi\,\rho_w\,g\,h}\Bigr)^{1/3}}\end{equation}
where $\delta$ is the thickness (m) and $h$ the suction (m). The film thins only as $h^{-1/3}$, so it survives at suctions where no pore is full.

\subhd{Corners.} In a pore of triangular or square cross-section, a meniscus of radius $r=\gamma/(\rho_wgh)$ fills each corner (the curvature is in one direction only, so the factor 2 of the cylinder is missing). The water in the corners forms continuous rivulets along the edges of the pores. The conductivity of the corners falls as $r^4$ and that of the films as $\delta^3$; together they explain why the measured conductivity of a dry soil lies above the value of the capillary model of Mualem, which counts only full pores.

\begin{kfig}{Cross-section of an angular pore in the dry range: water is held in the corners by menisci of radius $r=\gamma/(\rho_w g h)$ and as a thin film $\delta$ on the walls. Both carry flow along the pore after the centre has emptied.}
\draw[soil] (-3.4,-2.1) rectangle (3.4,2.5);
\fill[white] (-2.4,-1.4)--(2.4,-1.4)--(0,2.1)--cycle;\draw[ln] (-2.4,-1.4)--(2.4,-1.4)--(0,2.1)--cycle;
\fill[wtr] (-2.4,-1.4)--(-1.75,-1.4) to[bend left=30] (-2.05,-0.9)--cycle;\fill[wtr] (2.4,-1.4)--(1.75,-1.4) to[bend right=30] (2.05,-0.9)--cycle;\fill[wtr] (0,2.1)--(-0.33,1.62) to[bend right=30] (0.33,1.62)--cycle;
\draw[khcyan,line width=1.4pt] (-1.75,-1.38)--(1.75,-1.38);\draw[khcyan,line width=1.4pt] (-2.05,-0.9)--(-0.33,1.62);\draw[khcyan,line width=1.4pt] (2.05,-0.9)--(0.33,1.62);
\node[lab] at (0,-0.1) {air};
\draw[vec] (-1.6,-2.4)--(-2.05,-1.25);\node[lab,below] at (-1.6,-2.4) {corner meniscus $r$};
\draw[vec] (1.2,-2.4)--(0.6,-1.42);\node[lab,below] at (1.4,-2.4) {film $\delta$};
\end{kfig}

\exbox{Thickness of the water film}{Find the thickness of the adsorbed water film on a silicate surface ($A_H=6\times10^{-20}$ J) at suctions of 1 m, 10 m and 150 m (the last is the wilting point). Compare with the radius of the pore that is just full at 150 m.}\sol
\givenline{$A_H=6\times10^{-20}$ J ; $\rho_w=1000$ kg\,m$^{-3}$ ; $g=9.81$ m\,s$^{-2}$ ; $h=1$, $10$ and $150$ m}
\wkt\ the equilibrium thickness of the film (Equation~\ref{eq:pg-film}) is
\[\delta=\Bigl(\frac{A_H}{6\pi\,\rho_w\,g\,h}\Bigr)^{1/3}\ti{i}\]
The constant part of the denominator is \[6\pi\rho_wg=6\times3.1416\times1000\times9.81=1.849\times10^{5}\ \text{N\,m$^{-3}$}\]
\sh{1. Thickness at a suction of 1 m}
Substituting $h=1$ m in equation \rfi{i},
\[\delta=\Bigl(\frac{6\times10^{-20}}{1.849\times10^{5}\times1}\Bigr)^{1/3}=(3.245\times10^{-25})^{1/3}=6.87\times10^{-9}\ \text{m}=6.9\ \text{nm}\Ans\]
\sh{2. Thickness at a suction of 10 m}
\[\delta=\Bigl(\frac{6\times10^{-20}}{1.849\times10^{6}}\Bigr)^{1/3}=(3.245\times10^{-26})^{1/3}=3.19\times10^{-9}\ \text{m}=3.2\ \text{nm}\Ans\]
\sh{3. Thickness at the wilting point (150 m)}
\[\delta=\Bigl(\frac{6\times10^{-20}}{2.774\times10^{7}}\Bigr)^{1/3}=(2.163\times10^{-27})^{1/3}=1.29\times10^{-9}\ \text{m}=1.3\ \text{nm}\Ans\]
\sh{4. Radius of the pore just full at 150 m}
From the Young--Laplace law,
\[r=\frac{2\gamma}{\rho_w\,g\,h}=\frac{2\times0.0728}{1000\times9.81\times150}=9.9\times10^{-8}\ \text{m}=0.099\ \mu\text{m}\Ans\]
\[\therefore\quad\frac{r}{\delta}=\frac{99\ \text{nm}}{1.3\ \text{nm}}\approx76\]
The film at the wilting point is only four or five molecules of water thick (one molecule is about 0.28 nm), yet it remains connected over the whole surface of the soil and carries the slow flow of the dry range.

\section{The pore network as a lattice: percolation and critical paths}\label{sec:pg-perc}
The bundle of capillaries (Section~\ref{sec:bundle}) treats every pore as a straight tube that crosses the sample. A better picture is a \kw{network}: pore bodies at the nodes of a lattice, joined by throats along its bonds. Each node has $z$ neighbours, the \kw{coordination number}, which is about 4 to 8 in soils.

\subhd{Percolation threshold.} Let a fraction $p$ of the throats be open (filled with water, or with air for the gas phase). A connected path across the sample exists only when $p$ exceeds the \kw{percolation threshold} $p_c$. For a tree-like (Bethe) network $p_c=1/(z-1)$; for the simple cubic lattice ($z=6$) the threshold of bond percolation is $p_c=0.249$. Near the threshold the conductance of the network rises as a power law,
\begin{equation}\label{eq:pg-perc}\boxed{K\propto(p-p_c)^{t},\qquad t\approx2.0\ \text{in three dimensions}}\end{equation}
which is why the conductivity of a drying soil does not fall smoothly but collapses when the water-filled throats become rare.

\subhd{Critical path.} Because the conductance of a throat varies as $r^4$, the flow through a network with a wide range of sizes is controlled by the \kw{critical throat}: the smallest throat that must be used on the best connected path. Removing all throats narrower than $r_c$ just disconnects the network. The saturated conductivity is then set by $r_c$, not by the mean pore size; this is the idea behind the critical-path estimate $k\approx r_c^2/c$ with $c$ between 50 and 230 for granular media.

\begin{kfig}{A pore network on a square lattice: bodies at the nodes, throats on the bonds. Open throats (thick) form a spanning path from left to right only when their fraction exceeds the percolation threshold $p_c$; the narrowest throat on the best path (magenta) is the critical throat.}
\foreach \x in {0,...,6}{\foreach \y in {0,...,4}{\fill[khcyan!40] (\x*1.1,\y*0.9) circle (0.12);}}
\foreach \x in {0,...,5}{\foreach \y in {0,...,4}{\draw[thin ln] (\x*1.1+0.12,\y*0.9)--(\x*1.1+0.98,\y*0.9);}}
\foreach \x in {0,...,6}{\foreach \y in {0,...,3}{\draw[thin ln] (\x*1.1,\y*0.9+0.12)--(\x*1.1,\y*0.9+0.78);}}
\draw[khcyan!70!black,line width=2.2pt] (0,1.8)--(1.1,1.8)--(1.1,2.7)--(2.2,2.7)--(3.3,2.7)--(3.3,1.8);
\draw[khmag,line width=1.0pt] (3.3,1.8)--(4.4,1.8);
\draw[khcyan!70!black,line width=2.2pt] (4.4,1.8)--(4.4,0.9)--(5.5,0.9)--(6.6,0.9);
\node[lab,above] at (3.85,1.85) {$r_c$};
\draw[vec] (-0.9,1.8)--(-0.2,1.8);\draw[vec] (6.8,0.9)--(7.5,0.9);
\node[lab,below] at (3.3,-0.35) {open fraction $p>p_c$: a spanning path exists};
\end{kfig}

\exbox{Conductance near the percolation threshold}{In a network on a simple cubic lattice ($p_c=0.249$, exponent $t=2$) half of the throats are water-filled at a suction of 100 cm and 35 per cent at 300 cm. By what factor has the network conductance fallen? What fraction must remain filled for the water phase to stay connected?}\sol
\givenline{$p_c=0.249$ ; $t=2$ ; $p_1=0.50$ (at 100 cm) ; $p_2=0.35$ (at 300 cm)}
\wkt\ near the percolation threshold the conductance of the network is given by Equation~\ref{eq:pg-perc}:
\[K=C\,(p-p_c)^{t}\ti{i}\]
where $C$ is a constant of the network. Writing equation \rfi{i} for the two states and dividing, the constant cancels:
\[\frac{K_2}{K_1}=\Bigl(\frac{p_2-p_c}{p_1-p_c}\Bigr)^{t}\ti{ii}\]
Substituting the given values in equation \rfi{ii},
\[\frac{K_2}{K_1}=\Bigl(\frac{0.35-0.249}{0.50-0.249}\Bigr)^{2}=\Bigl(\frac{0.101}{0.251}\Bigr)^{2}=(0.4024)^2=0.162\Ans\]
\[\therefore\quad\frac{K_1}{K_2}=\frac1{0.162}=6.2\]
\sh{Fraction that keeps the water phase connected}
A spanning path of water exists only while $p>p_c$, i.e. while more than
\[p_c=0.249\quad\text{(about one throat in four)}\Ans\]
of the throats remain filled. The conductance has fallen by a factor of 6.2 although the filled fraction fell only from 0.50 to 0.35, which is why the conductivity of a drying soil collapses rather than falls smoothly.

\section{Evaporation from the pores: the two stages and the characteristic length}\label{sec:pg-evap}
A wet bare soil first loses water at nearly the rate of open water (\kw{stage 1}), limited by the energy at the surface. The surface stays wet because capillary flow lifts water from the large pores below through the small pores that remain full at the surface. When the drying front has receded to a depth at which the capillary pull of the small pores can no longer balance gravity and the viscous loss, the liquid connection breaks; evaporation then takes place below the surface and the vapour must diffuse up through a dry layer (\kw{stage 2}), at a rate that falls with the thickness of that layer.

The depth of the front at the end of stage 1 is the \kw{characteristic length} of the soil. If $r_s$ and $r_l$ are the radii of the small and large pores that form the connection, the capillary pressure difference that can hold the liquid against gravity gives
\begin{equation}\label{eq:pg-lg}\boxed{L_G=\frac{2\gamma}{\rho_w g}\Bigl(\frac{1}{r_s}-\frac{1}{r_l}\Bigr),\qquad L_C=\frac{L_G}{1+e_0/K_{\mathrm{eff}}}}\end{equation}
where $L_G$ is the length limited by gravity alone, $e_0$ the evaporation rate in stage 1 and $K_{\mathrm{eff}}$ the effective conductivity of the connecting path. A coarse sand with a narrow range of pores has a short characteristic length (a few centimetres) and passes quickly to stage 2; a fine-textured soil with a wide range of pores keeps stage 1 for days.

\begin{kfig}{The two stages of evaporation from a bare soil. Left: in stage 1 the small pores (narrow tubes) remain full up to the surface and feed it by capillarity from the large pores that are emptying. Right: the cumulative evaporation is linear in time in stage 1 and proportional to $\sqrt t$ in stage 2.}
\draw[soil] (0,0) rectangle (3.2,3.4);\draw[ln] (0,3.4)--(3.2,3.4);
\foreach \x in {0.5,1.4,2.3}{\fill[white] (\x,0) rectangle (\x+0.5,3.4);\draw[thin ln] (\x,0)--(\x,3.4);\draw[thin ln] (\x+0.5,0)--(\x+0.5,3.4);\fill[khcyan!25] (\x,0) rectangle (\x+0.5,1.6);}
\foreach \x in {1.1,2.0,2.9}{\fill[white] (\x,0) rectangle (\x+0.12,3.4);\fill[khcyan!40] (\x,0) rectangle (\x+0.12,3.4);}
\draw[gd] (0,1.6)--(3.2,1.6);\node[lab,left] at (0,1.6) {front};
\draw[<->,ln] (3.5,1.6)--(3.5,3.4);\node[lab,right] at (3.5,2.5) {$\le L_C$};
\foreach \x in {1.16,2.06,2.96}{\draw[vec] (\x,3.45)--(\x,4.05);}
\node[lab] at (1.6,4.3) {stage 1: evaporation at the surface};
\begin{scope}[xshift=5.6cm]
\draw[axis] (0,0)--(4.2,0);\draw[axis] (0,0)--(0,3.6);
\node[lab,below] at (2.1,-0.15) {time $t$};\node[lab,rotate=90] at (-0.4,1.8) {cumulative evaporation};
\draw[crv] (0,0)--(1.4,1.9);\draw[crv,smooth] plot coordinates {(1.4,1.9) (1.8,2.25) (2.3,2.55) (2.9,2.82) (3.6,3.08) (4.0,3.2)};
\draw[gd] (1.4,0)--(1.4,1.9);\node[lab,below] at (1.4,0) {$t_1$};
\node[lab,rotate=52] at (0.55,1.05) {stage 1};\node[lab] at (3.0,2.45) {stage 2, $\propto\sqrt t$};
\end{scope}
\end{kfig}

\exbox{Characteristic length of a sandy loam}{In a sandy loam the connection to the surface is made by pores of radius $r_s=20\ \mu$m, while the bulk of the water drains from pores of $r_l=200\ \mu$m. Find the length $L_G$ limited by gravity. If the stage-1 rate is $e_0=5$ mm\,d$^{-1}$ and the effective conductivity of the path $K_{\mathrm{eff}}=20$ mm\,d$^{-1}$, find $L_C$ and the water that can evaporate in stage 1 if the large pores hold 0.15 m$^3$\,m$^{-3}$.}\sol
\givenline{\[r_s=20\ \mu\text{m}=2\times10^{-5}\ \text{m}\]
\[r_l=200\ \mu\text{m}=2\times10^{-4}\ \text{m}\]
$e_0=5$ mm\,d$^{-1}$ ; $K_{\mathrm{eff}}=20$ mm\,d$^{-1}$ ; drainable porosity of the large pores $=0.15$}
\sh{Length limited by gravity}
\wkt\ the characteristic length limited by gravity alone is (Equation~\ref{eq:pg-lg})
\[L_G=\frac{2\gamma}{\rho_w\,g}\Bigl(\frac1{r_s}-\frac1{r_l}\Bigr)\ti{i}\]
The capillary length of water is
\[\frac{2\gamma}{\rho_w\,g}=\frac{2\times0.0728}{1000\times9.81}=1.484\times10^{-5}\ \text{m}^2\]
and the difference of the curvatures is
\[\frac1{r_s}-\frac1{r_l}=\frac1{2\times10^{-5}}-\frac1{2\times10^{-4}}=50\,000-5000=45\,000\ \text{m}^{-1}\]
Substituting in equation \rfi{i},
\[L_G=1.484\times10^{-5}\times45\,000=0.668\ \text{m}\Ans\]
\sh{Characteristic length with the viscous loss}
\[L_C=\frac{L_G}{1+e_0/K_{\mathrm{eff}}}=\frac{0.668}{1+5/20}=\frac{0.668}{1.25}=0.534\ \text{m}\Ans\]
\sh{Water that can evaporate in stage 1}
The large pores above the front release their water, so
\[\text{Water}=0.15\times L_C=0.15\times534\ \text{mm}=80\ \text{mm}\Ans\]
\[\therefore\quad\text{Duration of stage 1}=\frac{80}{5}=16\ \text{days}\]
A medium-textured soil can therefore lose a large part of a rain to evaporation before its surface dries, and a mulch, which cuts $e_0$, saves most in such soils.

\section{Diffusion of gases and vapour through the air-filled pores}\label{sec:pg-diff}
Oxygen reaches the roots, carbon dioxide and nitrous oxide leave the soil and water vapour crosses the dry layer of stage 2 by \kw{molecular diffusion} through the air-filled pores. Fick's law for the soil is written with an effective coefficient $D_s$:
\begin{equation}\label{eq:pg-fick}\boxed{J=-D_s\,\frac{\mathrm dC}{\mathrm dz},\qquad \frac{D_s}{D_0}=\frac{\varepsilon^{10/3}}{\phi^{2}}\ \ \text{(Millington and Quirk)}}\end{equation}
where $J$ is the flux (g\,m$^{-2}$\,s$^{-1}$), $C$ the concentration in the soil air, $D_0$ the diffusion coefficient in free air, $\varepsilon=\phi-\theta$ the air-filled porosity and $\phi$ the total porosity. The ratio $D_s/D_0$ is the \kw{relative diffusivity}; it is the gas-phase counterpart of the relative conductivity and expresses, in one number, the reduced cross-section and the tortuosity of the air-filled pores. The older relation of Penman, $D_s/D_0=0.66\,\varepsilon$, overestimates the diffusion in wet soils because it ignores the loss of connection.

\begin{kfig}{Diffusion through the air-filled pores: the path of a gas molecule (magenta) winds around the solids (brown) and the water (blue); the effective diffusivity falls steeply as the air-filled porosity $\varepsilon$ falls.}
\draw[soil] (0,0) rectangle (6.4,3.4);
\foreach \x/\y/\r in {0.7/0.8/0.55,1.9/2.4/0.6,2.9/0.9/0.5,4.0/2.5/0.62,5.2/1.1/0.58,5.7/2.7/0.4,1.2/1.8/0.3,3.5/1.7/0.3}{\fill[white] (\x,\y) circle (\r);\fill[khcyan!30] (\x,\y) circle (\r*0.55);}
\draw[crv,smooth] plot coordinates {(0.2,3.3) (0.8,2.6) (1.2,1.9) (1.5,1.3) (2.3,1.4) (2.9,1.5) (3.4,2.1) (3.6,1.4) (4.3,1.6) (4.7,2.1) (5.2,1.7) (5.8,1.9) (6.3,0.3)};
\node[lab,above] at (0.2,3.4) {$C_1$};\node[lab,below] at (6.3,0) {$C_2<C_1$};
\node[lab,right] at (6.5,1.7) {$J=-D_s\,\mathrm dC/\mathrm dz$};
\end{kfig}

\exbox{Oxygen diffusion in a wet and a dry loam}{A loam has porosity $\phi=0.45$. Find the relative diffusivity by Millington and Quirk and by Penman at an air-filled porosity of 0.20, and the effective diffusivity of oxygen ($D_0=0.205$ cm$^2$\,s$^{-1}$). Repeat for $\varepsilon=0.10$, after rain.}\sol
\givenline{$\phi=0.45$ ; $D_0=0.205$ cm$^2$\,s$^{-1}$ ; $\varepsilon=0.20$ and $0.10$}
\wkt\ the relative diffusivity by Millington and Quirk (Equation~\ref{eq:pg-fick}) is
\[\frac{D_s}{D_0}=\frac{\varepsilon^{10/3}}{\phi^{2}}\ti{i}\]
and by Penman
\[\frac{D_s}{D_0}=0.66\,\varepsilon\ti{ii}\]
\sh{1. Air-filled porosity of 0.20}
From equation \rfi{i}, with $\phi^2=0.45^2=0.2025$,
\[\frac{D_s}{D_0}=\frac{0.20^{10/3}}{0.2025}=\frac{0.004678}{0.2025}=0.0231\Ans\]
\[\therefore\quad D_s=0.0231\times0.205=4.74\times10^{-3}\ \text{cm}^2\,\text{s}^{-1}\Ans\]
From equation \rfi{ii},
\[\frac{D_s}{D_0}=0.66\times0.20=0.132\Ans\]
which is $0.132/0.0231=5.7$ times the value of Millington and Quirk.
\sh{2. Air-filled porosity of 0.10 (after rain)}
\[\frac{D_s}{D_0}=\frac{0.10^{10/3}}{0.2025}=\frac{4.642\times10^{-4}}{0.2025}=2.29\times10^{-3}\Ans\]
\[\therefore\quad\frac{D_s(0.20)}{D_s(0.10)}=\frac{0.0231}{0.00229}=10.1\]
Halving the air-filled porosity cuts the diffusion of oxygen tenfold. This steep fall is why a waterlogged root zone becomes anaerobic within a day or two, and why denitrification (Chapter~\ref{chap:17}) is concentrated in the wettest periods.

